\documentclass[10pt,a4paper]{amsart}
\usepackage[margin=19mm]{geometry}
\usepackage{amsmath,amssymb,mathtools}
\usepackage[scr=rsfs,cal=euler]{mathalfa}
\usepackage{xfrac}
\usepackage[unicode,hidelinks,bookmarks=false]{hyperref}
\newcommand{\ie}{i.e.\ }
\newcommand{\cf}{cf.\ }
\newcommand{\resp}{respectively\ }
\newcommand{\Subsection}{\subsection}
\providecommand{\nobreakdash}{\nobreak\hbox{-}}
\setlength{\emergencystretch}{2em}
\multlinegap0pt
\usepackage{bm}
\usepackage{bbm}
\usepackage{dsfont}
\usepackage{xspace}
\usepackage{cancel}
\usepackage{mathtools}
\usepackage{amsthm}
\makeatletter
\@ifundefined{thm}{\newtheorem{thm}{Theorem}}{}
\@ifundefined{lemma}{\newtheorem{lemma}[thm]{Lemma}}{}
\makeatother
\newcommand{\bR}{\mathbb{R}}
\newcommand{\rar}{\rightarrow}
\title[Short spherical $t$-design curves]{Short spherical $t$-design curves}
\author{Emily J. King, Dustin G. Mixon}
\date{Source: arXiv:2607.15386v1 (CC BY 4.0)}
\begin{document}
\maketitle

\noindent\emph{Source and licence.} Short spherical $t$-design curves. Version arXiv:2607.15386v1 (CC BY 4.0), distributed under the Creative Commons Attribution 4.0 International licence, as recorded in the arXiv licence metadata for this exact version and shown on the abstract page https://arxiv.org/abs/2607.15386v1. Full rights notice: licenses/arxiv-2607.15386-CC-BY-4.0.txt.\par\smallskip

\noindent\textit{Editorial scope.} Counted unit: the Thm whose source label is \texttt{thm.sine ansatz} in the article (source0.tex lines 946--966, source label \texttt{thm.sine ansatz}), delivered together with the article's own complete original proof of it (source0.tex lines 968--1051, one \texttt{proof} environment of 84 lines). No mathematics is written, repaired or re-derived: statement and proof are the article's own text line for line. Required context retained uncounted: lem.ansatz (lines 850--882); eq.2 design condition (lines 859--863). Every definition, display and label of those lines is retained unchanged. Omitted from this package and not redistributed: 1--849, 864--945, 967--967, 1052--4154. Nothing inside a retained excerpt is omitted or altered: every retained body is delivered exactly as the article's lines are written, so no omission receipt of any kind accompanies this package. Citations: the retained excerpts carry no citation command at all, so no bibliography entry of the article is transcribed and no reference is redistributed. Reference adaptation: none was needed and none was made; every reference inside the delivered unit resolves inside it, so no unresolved reference or citation is produced. Printed slips kept literal: no printed word, symbol, hypothesis or number of the counted statement or of its proof was altered. Layout and class: the article's own class is \texttt{article} with packages including babel, geometry, amsmath, amsthm, amssymb, xcolor, hyperref, tikz, pgfplots; the wrapper is the lane's pinned \texttt{amsart} editorial wrapper at 10pt on A4, which restates the theorem-like environments the retained text needs and adds the article's own macro definitions that the retained text needs (\texttt{\textbackslash bR}, \texttt{\textbackslash rar}), with extra wrapper packages (bm, bbm, dsfont, xspace, cancel, mathtools). The delivered numbering is the unit's own. Assets: no retained span contains a figure environment, a graphics-inclusion command, a PDF-page inclusion, a table or a TikZ picture; no image file is copied into this package, the package declares no asset dependency and no image file is redistributed with it. Licence: version arXiv:2607.15386v1 of this article is distributed under the Creative Commons Attribution 4.0 International licence, as recorded in the arXiv licence metadata for this exact version and shown on the abstract page https://arxiv.org/abs/2607.15386v1. A full rights notice is delivered at licenses/arxiv-2607.15386-CC-BY-4.0.txt. Characters: every retained excerpt is pure ASCII, so the delivered engine, XeLaTeX with its default Latin Modern text font, reports no missing character.
\begin{lemma}
\label{lem.ansatz}
Given a positive integer $k$, consider any $\frac{\pi}{k+1}$-periodic $x_1,\ldots,x_k\in C^1(\mathbb{R})$ and any $\frac{\pi}{k+1}$-antiperiodic\footnote{A function $f\colon \bR \rar \bR$ is said to be \emph{$p$-antiperiodic} if $f(x+p)=-f(x)$ for all $x \in \bR$.} $z\in C^1(\mathbb{R})$ such that
\begin{equation}
\label{eq.norm condition}
\sum_{j=1}^kx_j(\theta)^2+z(\theta)^2
=1
\end{equation}
for every $\theta\in\mathbb{R}$ and
\begin{equation}
\label{eq.2 design condition}
\int_0^{2\pi}\bigg(x_j(\theta)^2-\frac{2}{2k+1}\bigg)\sqrt{\sum_{\ell=1}^k \big(x_\ell'(\theta)^2+\ell^2x_\ell(\theta)^2\big)+z'(\theta)^2}\,d\theta
=0
\end{equation}
for each $j\in\{1,\ldots,k\}$.
Then $\gamma\colon[0,2\pi)\to S^{2k}$ defined by
\begin{align*}
\gamma(\theta)
\,&:=\,\big(
\,x_1(\theta)\cos\theta,
~x_1(\theta)\sin\theta,
~x_2(\theta)\cos2\theta,
~x_2(\theta)\sin2\theta,
~\cdots\\
&\hspace{1.6in}
~\cdots,
~x_k(\theta)\cos k\theta,
~x_k(\theta)\sin k\theta,
~z(\theta)
\,\big)
\end{align*}
is a $2$-design curve.
\end{lemma}

\begin{thm}
\label{thm.sine ansatz}
For every positive integer $k$, there exists $a_k\in(0,1)$ such that
\[
x_1(\theta)
=\cdots
=x_k(\theta)
=\sqrt{\frac{1-z(\theta)^2}{k}},
\qquad
z(\theta)=a_k\sin(k+1)\theta
\]
satisfy the hypotheses of Lemma~\ref{lem.ansatz}.
In particular, 
\[
a_1^2\in(0.820,0.821)
\qquad\text{and}\qquad
a_k^2\in\big(\tfrac{1}{k}-\tfrac{1}{2k^3},\tfrac{1}{k}+\tfrac{1}{4k^3}\big)
\quad
\forall\,k\geq2.
\]
\end{thm}

\begin{proof}
Note that $x_1,\ldots,x_k$, and $z$ satisfy all of the hypotheses of Lemma~\ref{lem.ansatz} except possibly~\eqref{eq.2 design condition}.
For this final constraint, we put $q:=a^2$ and change variables $u:=(k+1)\theta$ so that it suffices for the following quantity to vanish:
\[
H_k(q)
:=
\int_0^{2\pi}
\big(
\tfrac{1-q\sin^2u}{k}
-\tfrac{2}{2k+1}
\big)
\sqrt{
\tfrac{(k+1)(2k+1)}{6}\big(1-q\sin^2u\big)
+\tfrac{q(k+1)^2\cos^2u}{1-q\sin^2u}
}\,du.
\]
When $k=1$, one may use interval arithmetic to obtain
\[
H_1(0.820)>0
\qquad\text{and}\qquad
H_1(0.821)<0,
\]
and so the desired bounds on $a_1^2$ follow from the intermediate value theorem.

Now assume $k\geq 2$.
Again by the intermediate value theorem, it suffices to show
\[
H_k\bigg(\frac{1}{k}-\frac{1}{2k^3}\bigg)>0
\qquad\text{and}\qquad
H_k\bigg(\frac{1}{k}+\frac{1}{4k^3}\bigg)<0.
\]
To this end, we parameterize the intervening interval by
\[
q=\frac{1}{k}+\frac{c}{k^3},
\qquad
-\frac{1}{2}\leq c\leq \frac{1}{4}.
\]
We approximate the integrand of $H_k(\frac{1}{k}+\frac{c}{k^3})$ with a function that depends polynomially on $x:=1/k$.
The Taylor expansion of the integrand $F$ of $H_k(x+cx^3)$ about $x=0$ takes the form
\[
F(x,c,u)
=
C_1(u)\cdot x+C_2(u)\cdot x^2+C_3(c,u)\cdot x^3+R(x,c,u),
\]
and then integrating over $u$ gives
\begin{equation}
\label{eq.H expansion}
H_k\bigg(\frac1k+\frac{c}{k^3}\bigg)
=
-\frac{\sqrt3\pi(16c+1)}{48k^3}
+\underbrace{\int_0^{2\pi}
R(x,c,u)
\,du
}_{E(x,c)},
\qquad
x=\frac{1}{k}.
\end{equation}
One may use interval arithmetic to obtain
\[
\bigg|\frac{E(x,c)}{x^4}\bigg|
<1,
\qquad
0<x\leq\frac{1}{2},
\qquad
-\frac{1}{2}\leq c\leq\frac{1}{4}.
\]
Taking $c=-\frac{1}{2}$ and $c=\frac{1}{4}$ in~\eqref{eq.H expansion} and applying this bound, we get
\[
H_k\bigg(\frac{1}{k}-\frac{1}{2k^3}\bigg)
=\frac{7\sqrt{3}\pi}{48k^3}
+E\bigg(\frac{1}{k},-\frac{1}{2}\bigg)
>\frac{7\sqrt{3}\pi}{48k^3}
-\frac{1}{k^4}
>0,
\]
\[
H_k\bigg(\frac{1}{k}+\frac{1}{4k^3}\bigg)
=-\frac{5\sqrt{3}\pi}{48k^3}
+E\bigg(\frac{1}{k},\frac{1}{4}\bigg)
<-\frac{5\sqrt{3}\pi}{48k^3}+\frac{1}{k^4}
<0,
\]
where the inequalities use the fact that $k\geq 2$.
\end{proof}

\end{document}
